\begin{figure}[htbp]
\centering
\begin{tikzpicture}[
  box/.style={draw=NimbusBlue,fill=NimbusLight,rounded corners=2pt,
    text width=4.9cm,minimum height=17mm,align=center,font=\small},
  arrow/.style={-{Latex[length=2mm]},draw=NimbusBlue,thick}]
  \node[box] (a) at (0,0) {GPU A: owned horizontal tile\\complete vertical columns};
  \node[box] (b) at (7.2,0) {GPU B: owned horizontal tile\\complete vertical columns};
  \draw[arrow] ([yshift=3mm]a.east)--node[above,font=\scriptsize]{stage $s$ halo}([yshift=3mm]b.west);
  \draw[arrow] ([yshift=-3mm]b.west)--node[below,font=\scriptsize]{stage $s$ halo}([yshift=-3mm]a.east);
  \node[box] (stage) at (3.6,-2.4) {Validate the shared stage\\global timestep and owned budgets};
  \draw[arrow] (a.south)--(stage.north west);
  \draw[arrow] (b.south)--(stage.north east);
  \node[box,below=8mm of stage] (commit) {One accepted endpoint\\consistent output and checkpoint};
  \draw[arrow] (stage)--(commit);
\end{tikzpicture}
\caption{Proposed two-tile solver decomposition. Halo exchange can use a
validated direct path or explicit host staging. Arrows describe data
dependencies, not a mandatory global barrier after every local operator.
Each tile owns its physical cells once; copied halos carry no additional
atmospheric inventory. This is a design diagram, not a current Nimbus feature.}
\label{fig:architecture-multigpu}
\end{figure}
